% =====================================================================
\chapter{Información y entropía}
\label{tema:1}
% =====================================================================

Este tema define la medida de la información introducida por Shannon~\cite{shannon} y sus versiones para varias variables: entropía, entropía conjunta y entropía condicionada. Se sigue~\cite[§2.1, §2.2 y §2.5]{ct}; la interpretación de la entropía como número de preguntas y la discusión del contenido de información proceden de~\cite[cap.~2 y~4]{mackay} y de la primera clase de~\cite{camit}. En todo el tema, $X$ es una variable aleatoria discreta con valores en un alfabeto finito $\X$ y función de masa $p(x)=P(X=x)$.

\section{El contenido de información de un suceso}

Un suceso seguro no informa de nada; un suceso improbable, al ocurrir, informa mucho. Se busca una función $h$ que asigne a cada probabilidad $p\in(0,1]$ una cantidad de información y que cumpla tres requisitos naturales:
\begin{enumerate}
  \item $h(1)=0$: lo que es seguro no aporta información;
  \item $h$ es decreciente: cuanto menos probable es un suceso, más informa;
  \item $h(pq)=h(p)+h(q)$: la información de dos sucesos independientes es la suma de las informaciones.
\end{enumerate}

\begin{proposicion}[Caracterización del contenido de información]\label{prop:caracterizacion}
Sea $h\colon(0,1]\to\R$ decreciente (en sentido amplio) y tal que $h(pq)=h(p)+h(q)$ para todo $p,q\in(0,1]$. Entonces existe una constante $c\ge 0$ tal que $h(p)=-c\log_2 p$ para todo $p$.
\end{proposicion}

\begin{proof}
Sea $g(t)=h(2^{-t})$ para $t\ge 0$. Como $2^{-(s+t)}=2^{-s}2^{-t}$, la hipótesis da $g(s+t)=g(s)+g(t)$, y $g$ es creciente porque $t\mapsto 2^{-t}$ es decreciente y $h$ también. Sea $c=g(1)$. Tomando $s=t=0$ se obtiene $g(0)=0$, y por inducción $g(n)=n\,c$ para todo entero $n\ge 0$. Si $m,n$ son enteros positivos, $n\,g(m/n)=g(m)=m\,c$, luego $g(r)=c\,r$ para todo racional $r\ge 0$; en particular $c=g(1)\ge g(0)=0$. Para $t\ge 0$ real y racionales $r_1\le t\le r_2$, la monotonía da $c\,r_1\le g(t)\le c\,r_2$; haciendo tender $r_1$ y $r_2$ a $t$ se obtiene $g(t)=c\,t$. Finalmente, para $p\in(0,1]$ se toma $t=-\log_2 p\ge 0$ y resulta $h(p)=g(t)=-c\log_2 p$.
\end{proof}

La constante $c$ solo fija la unidad de medida. Se adopta $c=1$.

\begin{definicion}[Contenido de información]\label{def:sorpresa}
El \emph{contenido de información} o \emph{sorpresa} del resultado $x$, con $p(x)>0$, es
\[
h(x)=-\log_2 p(x)\quad\text{bits}.
\]
Con el logaritmo natural la unidad es el \emph{nat}; como $\ln p=\ln 2\cdot\log_2 p$, se tiene $1\text{ nat}=1/\ln 2\approx 1{,}443$ bits.
\end{definicion}

\begin{center}
\begin{tabular}{@{}lcc@{}}
\toprule
Suceso & $p$ & $h$ (bits) \\
\midrule
cara en una moneda equilibrada & $1/2$ & $1$ \\
un~3 en un dado equilibrado & $1/6$ & $2{,}585$ \\
uno entre 8 resultados equiprobables & $1/8$ & $3$ \\
un token al que el modelo asigna $0{,}9$ & $0{,}9$ & $0{,}152$ \\
un token al que el modelo asigna $0{,}1$ & $0{,}1$ & $3{,}322$ \\
uno entre $1024$ resultados equiprobables & $2^{-10}$ & $10$ \\
\bottomrule
\end{tabular}
\end{center}

Si $p(x)=2^{-k}$, el resultado $x$ puede identificarse con $k$ preguntas de respuesta sí o no que dividan por la mitad, en cada paso, el conjunto de candidatos. El contenido de información generaliza esta cuenta a probabilidades arbitrarias~\cite[cap.~4]{mackay}.

\begin{observacion}[Conexión con los LLM]
La pérdida de un modelo de lenguaje en la posición $t$ es $-\ln q(x_t\mid x_{1:t-1})$ (módulo~1, tema~5): la sorpresa, en nats, que el modelo experimenta ante el token que realmente aparece. Entrenar un LLM consiste en reducir la sorpresa media ante el texto real.
\end{observacion}

\section{Entropía}

\begin{definicion}[Entropía]\label{def:entropia}
La \emph{entropía} de $X$ es la sorpresa esperada:
\[
H(X)=H(p)=\E\bigl[-\log p(X)\bigr]=-\sum_{x\in\X}p(x)\log p(x),
\]
con el convenio $0\log 0=0$, justificado por $\lim_{t\to 0^{+}}t\log t=0$.
\end{definicion}

La entropía depende solo de las probabilidades, no de los valores que toma $X$. Mide la incertidumbre sobre $X$ antes de observarla o, de forma equivalente, la información que se obtiene, en media, al observarla.

\begin{ejemplo}\label{ej:entropia-tres}
Si $p=(1/2,1/4,1/4)$, entonces $H(p)=\tfrac12\cdot 1+\tfrac14\cdot 2+\tfrac14\cdot 2=1{,}5$ bits. Un esquema de preguntas óptimo pregunta primero si $X$ es el primer resultado; si no lo es, hace una segunda pregunta. El número medio de preguntas es $\tfrac12\cdot 1+\tfrac12\cdot 2=1{,}5$, que coincide con $H(p)$.
\end{ejemplo}

\begin{ejemplo}[Entropía binaria]\label{ej:binaria}
Si $X\sim\mathrm{Ber}(\theta)$, su entropía es la \emph{función de entropía binaria}
\[
H_b(\theta)=-\theta\log\theta-(1-\theta)\log(1-\theta).
\]
Vale $H_b(0{,}1)=0{,}469$, $H_b(0{,}25)=0{,}811$ y $H_b(0{,}5)=1$ bit. La función es cóncava, simétrica respecto de $\theta=1/2$, se anula en $\theta\in\{0,1\}$ y alcanza su máximo, 1 bit, en $\theta=1/2$ (\cref{fig:binaria}).
\end{ejemplo}

\begin{figure}[h]
\centering
\begin{tikzpicture}[x=7cm,y=3.2cm,font=\small]
  \draw[->] (0,0) -- (1.06,0) node[right] {$\theta$};
  \draw[->] (0,0) -- (0,1.15) node[above] {$H_b(\theta)$ (bits)};
  \foreach \t in {0.25,0.5,0.75,1} \draw (\t,0.02) -- (\t,-0.02) node[below] {$\t$};
  \node[below left] at (0,0) {$0$};
  \draw (0.01,1) -- (-0.01,1) node[left] {$1$};
  \draw[dashed,regla] (0.5,0) -- (0.5,1) (0,1) -- (0.5,1);
  \draw[thick,domain=0.0005:0.9995,samples=160,smooth] plot (\x,{-(\x*ln(\x)+(1-\x)*ln(1-\x))/ln(2)});
  \fill (0.1,0.469) circle (0.6pt) node[right=2pt] {$0{,}469$};
  \fill (0.25,0.811) circle (0.6pt) node[right=2pt] {$0{,}811$};
\end{tikzpicture}
\caption{Función de entropía binaria del \cref{ej:binaria}. Una moneda muy sesgada es casi previsible y contiene poca información por lanzamiento.}
\label{fig:binaria}
\end{figure}

\begin{proposicion}[Cotas elementales]\label{prop:cotas-entropia}
Para toda distribución $p$ sobre un alfabeto de $K$ símbolos,
\begin{enumerate}
  \item $H(p)\ge 0$, con igualdad si y solo si $p$ es degenerada, es decir, $p(x_0)=1$ para algún $x_0$;
  \item $H(p)\le\log K$, con igualdad si y solo si $p$ es la distribución uniforme.
\end{enumerate}
\end{proposicion}

\begin{proof}
(i) Si $0<p(x)\le 1$, entonces $-\log p(x)\ge 0$, luego cada sumando $-p(x)\log p(x)$ es no negativo y $H(p)\ge 0$. La suma se anula si y solo si cada sumando se anula, es decir, si cada $p(x)$ vale~0 o~1; como las probabilidades suman~1, exactamente una de ellas vale~1. (ii) Es consecuencia de la desigualdad de Gibbs del \cref{tema:2}; la demostración se propone en el \cref{eje:2-gibbs-ln}.
\end{proof}

\begin{center}
\begin{tabular}{@{}lc@{}}
\toprule
Distribución & Entropía (bits) \\
\midrule
uniforme sobre 27 símbolos (26 letras y el espacio) & $\log 27=4{,}755$ \\
$(1/2,1/4,1/4)$ & $1{,}5$ \\
moneda con $P(\text{cara})=0{,}1$ & $0{,}469$ \\
distribución degenerada & $0$ \\
\bottomrule
\end{tabular}
\end{center}

\section{Entropía conjunta y entropía condicionada}

Sean $X$ e $Y$ variables discretas con función de masa conjunta $p(x,y)$, marginales $p(x)$ y $p(y)$ y condicionada $p(y\mid x)=p(x,y)/p(x)$ cuando $p(x)>0$.

\begin{definicion}[Entropía conjunta y condicionada]\label{def:conjunta}
La \emph{entropía conjunta} de $(X,Y)$ y la \emph{entropía de $Y$ condicionada a $X$} son
\begin{gather*}
H(X,Y)=-\sum_{x,y}p(x,y)\log p(x,y),\\
H(Y\mid X)=\sum_{x}p(x)\,H(Y\mid X=x)=-\sum_{x,y}p(x,y)\log p(y\mid x),
\end{gather*}
donde $H(Y\mid X=x)=-\sum_y p(y\mid x)\log p(y\mid x)$ es la entropía de la distribución condicionada.
\end{definicion}

La entropía condicionada es la incertidumbre que queda sobre $Y$ cuando se conoce $X$, promediada sobre los valores de $X$.

\begin{teorema}[Regla de la cadena para la entropía]\label{teo:cadena-entropia}
$H(X,Y)=H(X)+H(Y\mid X)$.
\end{teorema}

\begin{proof}
Para cada par con $p(x,y)>0$ se tiene $p(x,y)=p(x)\,p(y\mid x)$, luego $\log p(x,y)=\log p(x)+\log p(y\mid x)$. Multiplicando por $-p(x,y)$ y sumando,
\[
H(X,Y)=-\sum_{x,y}p(x,y)\log p(x)-\sum_{x,y}p(x,y)\log p(y\mid x).
\]
En el primer sumando, $\sum_y p(x,y)=p(x)$, de modo que vale $-\sum_x p(x)\log p(x)=H(X)$. El segundo es $H(Y\mid X)$ por definición~\cite[§2.2]{ct}.
\end{proof}

\begin{corolario}[Regla de la cadena para secuencias]\label{cor:cadena-secuencias}
Para variables $X_1,\dots,X_n$,
\[
H(X_1,\dots,X_n)=\sum_{t=1}^{n}H(X_t\mid X_1,\dots,X_{t-1}).
\]
\end{corolario}

\begin{proof}
Por inducción sobre $n$, aplicando el \cref{teo:cadena-entropia} a $X=(X_1,\dots,X_{n-1})$ e $Y=X_n$.
\end{proof}

El \cref{cor:cadena-secuencias} es la versión informacional de la factorización autorregresiva del módulo~1, $P(x_{1:T})=\prod_t P(x_t\mid x_{1:t-1})$: tomando $-\log$ y esperanzas en esa factorización se obtiene exactamente esta suma. La incertidumbre de un texto completo es la suma de las incertidumbres de cada token dado su contexto.

\begin{ejemplo}\label{ej:dos-bits}
Sean $X$ e $Y$ dos bits que coinciden con probabilidad $0{,}8$: $p(0,0)=p(1,1)=0{,}4$ y $p(0,1)=p(1,0)=0{,}1$. Las marginales son uniformes, luego $H(X)=H(Y)=1$ bit. La entropía conjunta es
\[
H(X,Y)=-2\cdot 0{,}4\log 0{,}4-2\cdot 0{,}1\log 0{,}1=0{,}8\cdot 1{,}322+0{,}2\cdot 3{,}322=1{,}722\text{ bits}.
\]
Por el \cref{teo:cadena-entropia}, $H(Y\mid X)=1{,}722-1=0{,}722$ bits, que coincide con $H_b(0{,}2)$: conocido $X$, el valor de $Y$ es el mismo salvo con probabilidad $0{,}2$.
\end{ejemplo}

\begin{observacion}
En general $H(Y\mid X)\ne H(X\mid Y)$, pero $H(X)-H(X\mid Y)=H(Y)-H(Y\mid X)$, porque ambas diferencias valen $H(X)+H(Y)-H(X,Y)$ por la regla de la cadena aplicada en los dos órdenes. Esta cantidad común es la información mutua del \cref{tema:2}.
\end{observacion}

\section{Tasa de entropía y redundancia del lenguaje}

\begin{definicion}[Tasa de entropía]\label{def:tasa}
La \emph{tasa de entropía} de un proceso $X_1,X_2,\dots$ es
\[
\mathcal{H}=\lim_{n\to\infty}\frac1n H(X_1,\dots,X_n),
\]
cuando el límite existe. Se mide en bits por símbolo.
\end{definicion}

Si las $X_t$ son independientes e idénticamente distribuidas, $H(X_1,\dots,X_n)=nH(X_1)$ y $\mathcal{H}=H(X_1)$. En un texto, en cambio, los símbolos dependen fuertemente del contexto y $\mathcal{H}$ es mucho menor que la entropía de un símbolo aislado~\cite[cap.~4]{ct}. Shannon~\cite{shannon51} estimó la tasa de entropía del inglés pidiendo a personas que adivinaran un texto letra a letra, y obtuvo entre $0{,}6$ y $1{,}3$ bits por letra, muy por debajo de los $\log 27=4{,}755$ bits de un alfabeto de 27 símbolos equiprobables. La diferencia mide la \emph{redundancia} del lenguaje: la parte del texto que es previsible a partir de lo anterior.

\begin{observacion}[Conexión con los LLM]
La pérdida media por token de un modelo de lenguaje, medida sobre mucho texto, es una estimación de una entropía cruzada que, por la desigualdad de Gibbs del \cref{tema:2}, no puede ser menor que la tasa de entropía del texto. Un modelo que usa un contexto más largo puede aprovechar más dependencias y acercarse más a $\mathcal{H}$; este es el sentido informacional de la mejora de los modelos de $n$-gramas a los Transformers.
\end{observacion}

\begin{lecturas}
\item Cover y Thomas~\cite{ct}, §2.1 «Entropy», §2.2 «Joint Entropy and Conditional Entropy» y §2.5 «Chain Rules for Entropy, Relative Entropy, and Mutual Information».
\item MacKay~\cite{mackay}, cap.~2 «Probability, Entropy, and Inference» y cap.~4 «The Source Coding Theorem».
\item Universidad de Cambridge, \emph{Information Theory}~\cite{camit}, temas «Information, probability, uncertainty, and surprise» y «Entropy of discrete variables».
\item Shannon~\cite{shannon}, partes~I.6 y~I.7, y Shannon~\cite{shannon51}.
\end{lecturas}

\section*{Ejercicios del tema 1}
\addcontentsline{toc}{section}{Ejercicios del tema 1}

\begin{ejercicio}{1}{eje:1-aditiva}
Explique por qué la información de dos sucesos independientes debe ser la suma de sus informaciones y qué función transforma el producto de probabilidades en esa suma. ¿Qué papel desempeña la base del logaritmo?
\end{ejercicio}

\begin{ejercicio}{1}{eje:1-archivos}
Un archivo de $1$~MB contiene solo ceros y otro, del mismo tamaño, contiene bits aleatorios independientes y equiprobables. ¿Cuál de los dos puede comprimirse más sin pérdidas? Razónelo con la entropía por bit de cada fuente.
\end{ejercicio}

\begin{ejercicio}{1}{eje:1-cambridge}
Calcule, en bits, la entropía de las siguientes variables: (a) el nivel de gris de un píxel, uniforme entre los enteros $0$ y $255$; (b) la clasificación de un ser humano como mamífero o no mamífero; (c) el sexo de un insecto de una especie con tres sexos de probabilidades $1/4$, $1/4$ y $1/2$; (d) la clasificación de una persona según tenga o no más edad que la mediana de su población.
\fuente{adaptado de Universidad de Cambridge, \emph{Information Theory}~\cite{camit}, \emph{Examples Sheet}, \emph{Supervision}~1, ejercicio~3.}
\end{ejercicio}

\begin{ejercicio}{2}{eje:1-binaria}
(a) Calcule la entropía de una moneda con $P(\text{cara})=1/4$. (b) Derive $H_b(\theta)$ y demuestre que su único máximo en $(0,1)$ está en $\theta=1/2$.
\end{ejercicio}

\begin{ejercicio}{2}{eje:1-geometrica}
Sea $X$ uniforme en $\{1,\dots,8\}$ y sea $Y$, independiente de $X$, con $P(Y=k)=2^{-k}$ para $k=1,2,3,\dots$ (a) ¿Cuál de las dos variables tiene más incertidumbre? Calcule $H(X)$ y $H(Y)$. (b) Calcule $H(X,Y)$ y $H(Y\mid X)$.
\fuente{adaptado de Universidad de Cambridge, \emph{Information Theory}~\cite{camit}, \emph{Examples Sheet}, \emph{Supervision}~3, ejercicio~2 (la información mutua se trata en el \cref{tema:2}).}
\end{ejercicio}

\begin{ejercicio}{3}{eje:1-funcion}
Sea $g\colon\X\to\mathcal{Y}$ una función. Demuestre que $H(g(X))\le H(X)$, con igualdad si y solo si $g$ es inyectiva sobre el conjunto $\{x:p(x)>0\}$. Indicación: calcule $H(X,g(X))$ de dos maneras con la regla de la cadena. Interprete el resultado para un preprocesado del texto que convierte todas las letras en minúsculas.
\end{ejercicio}
